\begin{table}[!htbp]
\centering
\begin{threeparttable}
\caption{SME firm turnover Pre vs Post the \policyCutoffMonth{} cutoff, by class. Discriminating test between the selection and strict-invariance readings of $F_c^{\text{SME,Post}}$.}
\label{tab:pharma_firm_turnover}
\footnotesize
\setlength{\tabcolsep}{6pt}
\begin{tabular}{lcccc}
\toprule
 & Pre-period & Post-period & Continuing & New (Post only) \\
Class & SME firms & SME firms & (in both)  & (\% of Post firms / Post bids) \\
\midrule
Non-pharma & 313 & 520 & 290 & 230 (44.2\% / 23.0\% bids)\\
Pharma & 38 & 97 & 37 & 60 (61.9\% / 37.8\% bids)\\
\bottomrule
\end{tabular}
\begin{tablenotes}\footnotesize
\item Sample: Preg\~ao SME bidders in the \structuralWindowM-month structural window. Pre = $m < 698$ (Sep 2016--Feb 2018); Post = $m \ge 698$ (Mar 2018--Aug 2019). Continuing firms appear in both periods; New firms appear only Post; the difference is exit.
\item \textbf{Discriminating reading.} Under the equilibrium-selection interpretation of $F_c^{\text{SME,Post}}$ (Assumption~\ref{a:setaside}, main specification), the protected regime should induce previously unprofitable firms to enter, so a substantial share of post-period bids should come from new firms. Under the strict-invariance reading, post-period bids should be dominated by continuing firms with the same cost primitive. The pharma column compared to the non-pharma column is the within-paper discriminating test.
\end{tablenotes}
\end{threeparttable}
\end{table}
